{\color{red}\noindent[\textbf{Fix P3, P4} (fix\_poles\_dglap Steps 4, 7, 11, 12). In the notation of Eq.~\eqref{eq:fx-DGLAP}, $\frac{d^{3}N_{\rm RC}}{d^{3}k}=+\mathcal{C}\,\frac{\beta_{0}}{2}\,\mathcal{F}(1)$ exactly (same phase convention $e^{-ik\cdot(y-y')}$, same $\mu$, $2\gamma_{E}+\log\frac{r^{2}\mu^{2}}{4}=\log\frac{r^{2}\mu^{2}}{c_{0}^{2}}$), i.e.\ it is minus the $\delta(1-\xi)$ part of $N_{\rm DGLAP}$, added and subtracted:
\begin{align}
\frac{d^{3}N_{\rm DGLAP}}{d^{3}k}+\frac{d^{3}N_{\rm RC}}{d^{3}k}=-\mathcal{C}\int_{0}^{1}d\xi\Big\{\mathcal{P}_{gg+}(\xi)\,\mathcal{F}(\xi)-2N_{c}\frac{1-\xi}{\xi}\mathcal{F}(0)\Big\}\;{\color{green!50!black}\cancel{-\,\mathcal{C}\,\frac{\beta_{0}}{2}\,\mathcal{F}(1)}}\;{\color{green!50!black}\cancel{+\,\mathcal{C}\,\frac{\beta_{0}}{2}\,\mathcal{F}(1)}}\,,\qquad\mathcal{P}_{gg+}\equiv\mathcal{P}_{gg}-\frac{\beta_{0}}{2}\delta(1-\xi)\,.\notag
\end{align}
\textbf{What the ``running coupling'' term is} (answers to review Q2). (1) Its pole, as it arises in Group II Row IV, is collinear (IR) (the reviewer is right on this point). (2) What renormalises the coupling is the UV half of the (scaleless) normalisation of the dressed produced gluon, $Z_{3}$ of the produced gluon, consistent with the light-cone-gauge Ward identity $Z_{g}=Z_{3}^{-1/2}$; its IR half is the virtual $\frac{\beta_{0}}{2}\delta(1-\xi)$ which completes the real $\mathcal{P}_{gg+}$ to the $\overline{\rm MS}$ $\mathcal{P}_{gg}$ absorbed by the fragmentation counterterm. (3) The content of $N_{\rm RC}$ is the $\mu$ dependence: $\partial_{\log\mu^{2}}N_{\rm RC}=+\frac{\alpha_{s}\beta_{0}}{4\pi}\times$LO compensates the running of $g^{2}(\mu)$ in LO. In the classification of Kovner--Lublinsky--Skokov--Zhao (2308.15545): one $\beta_{0}$ (UV half of $Z_{3}$) is the $\beta$ function, the other (IR half, the $\frac{\beta_{0}}{2}\delta(1-\xi)$ of $\mathcal{P}_{gg}$) is DGLAP; here both are tied to the produced gluon and are equal in magnitude because the normalisation integral is scaleless. (4) The virtual $\frac{11}{3}$ terms of the one-loop wave function do \emph{not} give a running coupling at the emission vertex: they cancel between Group I Rows II and V (note at Group I Row II). Full pole table: \texttt{Divergent\_terms\_fixed}.]}\par
